\section{What Is Reproduced on Models, Exactly}
\label{sec:models}

Reproductions are worth having when they are exact and when the model is named. Three added by the audit are of that
kind, and each taught the program something about its own objects.

\textbf{The toric code from the mesh complex.} The periodic $D_4$ lattice has a natural cell complex: cells as
vertices, the twenty-four root directions as edges, and the triangles closed by pairs of roots whose sum is a root.
Qubits on edges, an $X$ check on every vertex star, a $Z$ check on every triangle. On the side-3 torus this is 972
qubits with every $X$ and $Z$ check overlapping on an even number of edges, and the rank computation over
$\mathbb{F}_2$ gives exactly four logical qubits, the first Betti number of the four-torus \cite{kitaev2003,
dennis2002}. An open patch has zero. The prediction was four at every side, and the measurement returned eight at
side 4. The explanation is the parity fact of Section~2: the even-sided lattice is two four-tori, each with its own
four. The code counted components before anyone had looked for them. The mesh, then, is a substrate on which
stabiliser codes live natively, and the hyperbolic versions \cite{breuckmann2016} are the next thing to build on a
closed quotient.

\textbf{The Ward identity and the Wilson lines.} Section~4 gave the numbers. What the reproduction taught is that the
zero-mode count of a lattice gauge operator on a torus is not one gauge mode per site less one, but that plus the
topology, $b_1 = 3$, which the naive prediction missed and the measurement supplied. A second experiment that tested only
that the zero-mode count was positive now tests it against $\text{sites} + 2$ exactly.

\textbf{The Ehrenfest theorem, with a packet that is a particle.} The walk's dispersion is $\cos E = \cos m \cos k$
with group velocity $v = \cos m \sin k / \sin E$, and the positive band is the eigenvector of the one-step unitary
with eigenvalue $e^{-iE}$. A Gaussian packet built from that band alone, at momentum $0.6$ and mass $0.5$ on 400
sites, is $99.9999\%$ positive-energy. The Ehrenfest prediction is computed from the packet's initial momentum
distribution, not read off the run: the mean momentum shifts by exactly $-F$ per step under a linear potential, and
the centroid advances by the distribution-averaged group velocity. Two facts had to be right for the prediction to
hold. A band packet sits at the Berry connection of its band, $-0.142$ cells here, measured and predicted alike, and
that offset moves as the force drifts the momentum. And the momentum kick lands between the coin and the shift, so
each step sees the midpoint momentum $k - F(t + \tfrac12)$: with the start-of-step momentum the residual is 0.19
cells and grows linearly with the force, with the midpoint it is 0.0024. Free, the centroid matches to
$9\times10^{-14}$ cells over 57 cells of travel. Under a force of 0.005 per step it matches to 0.0024 cells over 45,
with the momentum exact to $4\times10^{-16}$ and a Zener leak of $2.4\times10^{-5}$ into the other band. The controls
are the two ways to break the theorem: a force of 0.8 leaves 26\% of the weight in the band and the centroid 14
cells off, and the equal-chirality Gaussian ``at rest'', which the first attempt used and which followed no
classical path, is 21\% negative-energy and lags the classical trajectory by 19 cells. That seed was the reason the
first probe failed, and it is kept as the control that explains why \cite{ehrenfest1927}.

These are the shapes a good L2 has: an exact number, a control that breaks it, a prediction that was wrong once and
corrected by the measurement, and the model named in the title.
